\chapter*{Introduction}\addcontentsline{toc}{chapter}{Introduction}\label{sec:introduction}
The \gretlpy package is a software collection that
provides an interface to \lgretl \citep{gretl_manual}
for the \python{3.xx} (3.11, 3.12, 3.13, and 3.14 were tested)
language and consists of the following elements:
\begin{itemize}
	\item binary bindings written in \cpp,
	\item dedicated classes implemented in \cpp and exposed to \python
		(see Section~\ref{sec:class_py_model}),
	\item a set of utility functions implemented in \python
	(see Sections~\ref{sec:module_related_python}--\ref{sec:data_management_python}).
\end{itemize}
On \cpp level, we use \pkg{pyind11} \citep[see][for details]{pybind11}.

The package is distributed in \texttt{wheel} format
(\url{https://packaging.python.org/en/latest/glossary/#term-Wheel}).


\section*{Warning}
\begin{quote}
	This documentation is continuously updated but does
	not yet cover all \gretlpy~\gretlpyversion features.
\end{quote}


\section*{To do}
\begin{quote}
	\input{../todo.txt}
\end{quote}